% solutions/lem-fiber-root-degree-two.tex
% Standalone proof dossier for the ledger node lem:fiber-root-degree-two, which
% research/program/ledger.yaml carries with status: open.  Nothing in this file
% changes any ledger state.  This dossier concerns the fixed A_{m-1} root frame on
% the degree-<=2 polynomial quotient only; it is NOT about the all-frame
% gate q:conditional-fiber-frame.
\documentclass[../main.tex]{subfiles}
\begin{document}

% === SOLUTION HEADER =========================================================
% Certification is not recorded here: research/program/ledger.yaml owns the
% proofs[] mode and review path, and the review's front matter owns identities.
% No proofs[] record names this file, and that absence is the whole signal:
% this is a draft, unreviewed, and not yet a proof for ledger purposes.
%
% ledger-node : lem:fiber-root-degree-two
% refines : lem:fiber-root-degree-two   (modules/kls/45-conditional-fiber-frame.tex)
% bounded_by : none
% author : prover agent (Claude Fable 5), run id w4p04
% date : 2026-08-30
% ============================================================================
%
% Reader's notes below, deliberately not in field shape so that nothing here is
% parsed as header metadata.
%
%   No obstruction node fences this statement. The audit paragraph checks the
%   statement for consistency against prop:conditional-fiber-root-obstruction
%   and prop:sasada-negative-exchange.
%
%   This dossier complements prop:conditional-fiber-root-obstruction and feeds the
%   refutation-channel analysis of q:conditional-fiber-frame.
%
%   Dependencies: the theorem part is self-contained (Dirichlet moments and finite
%   group theory); the corollary part uses prop:conditional-fiber-root-obstruction
%   and lem:conditional-fiber-form through the certified dossier
%   solutions/conditional-fiber-frame-structure.tex.

\section*{Solution: the degree-two root-frame pencil identity on the isotropic uniform
simplex}
\label{sec:sol-lem-fiber-root-degree-two}

\paragraph{Scope.}
This dossier proves an exact finite-dimensional spectral identity for the $A_{m-1}$
root-frame conditional-fiber form of
Route~F (\ref{sec:conditional-fiber-frame}) restricted to polynomial test functions of
degree at most two, for \emph{every} $m\ge3$ simultaneously.  The statement was
suggested by exact rational computations at $m=3,\dots,15$ recorded in
provenance-stamped run artifacts (see Remark~\ref{rem:sol-frd2-artifacts}); no step of
any proof below uses those computations.  The main theorem is unconditional and
self-contained.  The corollary is stated \emph{only} as a conditional,
refutation-channel statement inside the candidate degree-$k$ dual-certificate framework
for Question~\ref{q:conditional-fiber-frame}; it makes no claim about the all-frame
gate itself, about degrees $k\ge3$, about non-polynomial tests, or about KLS.

\paragraph{Setting.}
Let $m\ge3$, let $\Delta_{m-1}=\{p\in\R_+^m:\sum_ip_i=1\}$, and let
$P=(P_1,\dots,P_m)$ be uniform on $\Delta_{m-1}$, i.e.\ $P\sim\Dir(1,\dots,1)$; write
$\mu$ for its law.  Put $H_0=\one^\perp$, $d=m-1$, $R_m=\sqrt{m(m+1)}$, and
\begin{equation}
  X=R_m\Bigl(P-\frac1m\one\Bigr)\in H_0,
  \label{eq:sol-frd2-isotropic}
\end{equation}
which is isotropic on $H_0$ (uniform-simplex root proposition of the certified dossier
\texttt{solutions/conditional-fiber-frame-structure.tex}; this normalization is not
used before the corollary).  The symmetric group $S_m$ acts on functions by permuting
the coordinates of $p$, and $\mu$ is exchangeable, so $S_m$ acts by isometries of
$L^2(\mu)$.

For $i<j$ let $\mathcal F_{ij}=\sigma\bigl((P_\ell)_{\ell\ne i,j}\bigr)$, let
$s_{ij}=P_i+P_j$ (an $\mathcal F_{ij}$-measurable variable, since
$s_{ij}=1-\sum_{\ell\ne i,j}P_\ell$), and write $\Var_{ij}$ and $\Cov_{ij}$ for
conditional variance and covariance given $\mathcal F_{ij}$.  Define, for
$f,g\in L^2(\mu)$ with all the conditional quantities finite,
\begin{equation}
  K(f,g)
  =\frac{12}{m^2(m+1)}\sum_{i<j}
    \E\,\frac{\Cov_{ij}(f,g)}{s_{ij}^{\,2}},
  \qquad
  G(f,g)=\Cov(f,g).
  \label{eq:sol-frd2-forms}
\end{equation}
The quadratic form $K(f,f)$ is exactly the root-frame conditional-fiber form
$\mathcal D_{\rm root}(f)$ of \eqref{eq:conditional-root-form}, as proved in the
certified dossier above (its equation defining the root form); for the main theorem,
\eqref{eq:sol-frd2-forms} is simply taken as the definition, so the theorem does not
depend on that identification.  Both $K$ and $G$ depend only on the $\mu$-equivalence
classes of $f,g$, vanish when one argument is a.e.\ constant, and are $S_m$-invariant
(exchangeability of $\mu$ and invariance of the family of pairs).

Let $W_{\rm poly}$ be the linear span of the monomials
$p_i$ ($1\le i\le m$) and $p_ip_j$ ($1\le i\le j\le m$), i.e.\ all polynomials of degree
at most $2$ with zero constant term, and define the \emph{degree-two quotient}
\begin{equation}
  V_{m,2}
  =\bigl\{[f]:f\in W_{\rm poly}\bigr\}
  \subset L^2(\mu)\big/\R\one,
  \label{eq:sol-frd2-V}
\end{equation}
where $[f]$ denotes the class of $f$ modulo additive constants.  On $V_{m,2}$ the form
$G=\Var$ is positive definite by construction, $K$ is positive semidefinite and finite
(Lemma~\ref{lem:sol-frd2-pair}), and both are $S_m$-invariant.  The generalized
minimum eigenvalue of the pencil is
\begin{equation}
  \lmin(K,G)\big|_{V_{m,2}}
  =\min_{0\ne v\in V_{m,2}}\frac{K(v,v)}{G(v,v)}.
  \label{eq:sol-frd2-lmin}
\end{equation}

\subsection*{Statements}

\begin{theorem}[degree-two root-frame pencil identity]
\label{thm:sol-fiber-root-degree-two}
For every $m\ge3$,
\begin{equation}
  \lmin(K,G)\big|_{V_{m,2}}
  =\frac{(m+2)(m+3)}{5m^2},
  \label{eq:sol-frd2-main}
\end{equation}
and the minimum is attained exactly on the one-dimensional line
$\R\cdot\bigl[\textstyle\sum_i p_i^2\bigr]
 =\R\cdot\bigl[\abs{X}^2-(m-1)\bigr]$
spanned by the radial quadratic.  In particular
$\lmin(K,G)|_{V_{m,2}}>\tfrac15$ for every $m\ge3$, and
$\lmin(K,G)|_{V_{m,2}}\to\tfrac15$ as $m\to\infty$.
\end{theorem}

The corollary below lives inside the candidate degree-$k$ dual-certificate framework
for the simplex falsification channel of Question~\ref{q:conditional-fiber-frame}
(the all-frame min--max $\Lambda_{m,k}$ and its exact dual certificates, recorded in
\texttt{research/explorations/2026-08-27-kls-route-prober-conditional-fiber-frame-w1f01.md},
eqs.~(42)--(49); that framework is itself a candidate object whose ledger admission is
pending).  We restate the needed objects to be self-contained.  Write $S(H_0)$ for the
unit sphere of $H_0$.  For $\theta\in S(H_0)$ and $f$ a polynomial, let
\begin{equation}
  \qJac_\theta[f]
  =\int_{\theta^\perp}
   \frac{\Var_{\mu_{\theta,z}}\bigl(f(z+T\theta)\bigr)}
        {\Var_{\mu_{\theta,z}}(T)}
   \dd\bar\mu_\theta(z)
  \label{eq:sol-frd2-qtheta}
\end{equation}
be the single-direction normalized fiber energy of the certified structural dossier
(with its null- and zero-variance-fiber conventions), a finite quantity for polynomial
$f$ by the certified factor-$4$ bound.  An \emph{admissible frame} is an even Borel
probability $\rho$ on $S(H_0)$ with $(m-1)\int\theta\theta^T\dd\rho=I_{H_0}$.  A
\emph{degree-$2$ dual certificate with objective $\eps$} is a finite family
$c_1,\dots,c_R\in V_{m,2}$, weights $w_1,\dots,w_R\ge0$, and a symmetric matrix
$M$ on $H_0$ such that
\begin{equation}
  \sum_{r=1}^R w_r\,G(c_r,c_r)=1,
  \qquad
  \sum_{r=1}^R w_r\,\qJac_\theta[c_r]\le\theta^TM\theta
  \quad\text{for all }\theta\in S(H_0),
  \qquad
  \Tr M\le\eps.
  \label{eq:sol-frd2-certificate}
\end{equation}
By the weak-duality computation reproduced in the proof below, such a certificate
forces every admissible frame $\rho$ (for which the averaged form is defined) to have
degree-two pencil gap at most $\eps$; a sequence of such certificates with
$\eps_m\to0$ at the fixed degree $k=2$ would therefore be the decisive fixed-degree
refuter of the route on the simplex.

\begin{corollary}[conditional: no degree-two dual refuter]
\label{cor:sol-frd2-no-degree-two-refuter}
Assume the identification, certified in
\texttt{solutions/conditional-fiber-frame-structure.tex} (nodes
\emph{lem:conditional-fiber-form} and
\emph{prop:conditional-fiber-root-obstruction}), of the pair-redistribution formula
\eqref{eq:sol-frd2-forms} with the conditional-fiber form of the admissible
$A_{m-1}$ root frame, i.e.
\begin{equation}
  K(f,f)=(m-1)\int_{S(H_0)}\qJac_\theta[f]\dd\rho_{\rm root}(\theta),
  \qquad
  \rho_{\rm root}=\frac1{m(m-1)}\sum_{i\ne j}\delta_{(e_i-e_j)/\sqrt2}.
  \label{eq:sol-frd2-root-identification}
\end{equation}
Then for every $m\ge3$:
\begin{enumerate}
\item every degree-$2$ dual certificate \eqref{eq:sol-frd2-certificate} has
  \begin{equation}
    \eps\;\ge\;\frac{(m+2)(m+3)}{5m^2}\;>\;\frac15;
  \end{equation}
\item consequently no sequence of degree-$2$ dual certificates with
  $\eps_m\to0$ exists, and within this framework any fixed-degree polynomial dual
  refutation of the conditional-fiber-frame route on the uniform simplex requires
  test degree $k\ge3$;
\item for every admissible frame $\rho$ whose averaged degree-two form
  $A_\rho(f,f)=(m-1)\int\qJac_\theta[f]\dd\rho(\theta)$ is defined on $V_{m,2}$, the
  all-frame degree-two min--max satisfies
  \begin{equation}
    \Lambda_{m,2}
    :=\sup_{\rho\ \text{admissible}}
      \ \min_{0\ne v\in V_{m,2}}\frac{A_\rho(v,v)}{G(v,v)}
    \;\ge\;\frac{(m+2)(m+3)}{5m^2}\;>\;\frac15 .
  \end{equation}
\end{enumerate}
This corollary asserts nothing about upper bounds on $\Lambda_{m,2}$, about
$\Lambda_{m,k}$ for $k\ge3$, about non-polynomial test functions, about the truth
value of Question~\ref{q:conditional-fiber-frame}, or about KLS.
\end{corollary}

\subsection*{Preliminaries: moments and pair fibers}

\begin{lemma}[Dirichlet monomial moments]
\label{lem:sol-frd2-moments}
For $a\in\mathbb N^m$ with $\abs a=\sum_ia_i$,
\begin{equation}
  \E\prod_{i=1}^mP_i^{a_i}
  =\frac{(m-1)!\,\prod_{i=1}^ma_i!}{(m-1+\abs a)!}.
  \label{eq:sol-frd2-moment-formula}
\end{equation}
In particular, with $D_r=m(m+1)\cdots(m+r-1)$ and distinct indices $i,j,k,l$:
\begin{equation}
\begin{gathered}
  \E P_i=\frac1m,\qquad
  \E P_i^2=\frac2{D_2},\qquad
  \E P_iP_j=\frac1{D_2},\\
  \E P_i^3=\frac6{D_3},\qquad
  \E P_i^2P_j=\frac2{D_3},\qquad
  \E P_i^4=\frac{24}{D_4},\qquad
  \E P_i^3P_j=\frac6{D_4},\\
  \E P_i^2P_j^2=\frac4{D_4},\qquad
  \E P_i^2P_jP_k=\frac2{D_4},\qquad
  \E P_iP_jP_kP_l=\frac1{D_4}.
\end{gathered}
  \label{eq:sol-frd2-moment-table}
\end{equation}
\end{lemma}

\begin{proof}
Let $E_1,\dots,E_m$ be i.i.d.\ standard exponentials and $S=\sum_iE_i$.  By the
classical Gamma--Dirichlet factorization, $(E_1/S,\dots,E_m/S)\sim\Dir(1,\dots,1)$ and
is independent of $S\sim\GammaLaw(m,1)$.  Hence
\[
  \prod_ia_i!
  =\E\prod_iE_i^{a_i}
  =\E\Bigl[S^{\abs a}\prod_iP_i^{a_i}\Bigr]
  =\E S^{\abs a}\cdot\E\prod_iP_i^{a_i}
  =\frac{(m-1+\abs a)!}{(m-1)!}\,\E\prod_iP_i^{a_i},
\]
using $\E S^{r}=\Gamma(m+r)/\Gamma(m)$.  Rearranging gives
\eqref{eq:sol-frd2-moment-formula}; the table is immediate.
\end{proof}

The following consequences are used repeatedly ($i,j,a,b$ distinct):
\begin{equation}
  \E s_{ij}=\frac2m,\qquad
  \E s_{ij}^2=\frac{2\cdot2+2\cdot1}{D_2}=\frac6{D_2},\qquad
  \E(P_a-P_b)^2=\frac{2\cdot2-2\cdot1}{D_2}=\frac2{D_2}.
  \label{eq:sol-frd2-small-moments}
\end{equation}

\begin{lemma}[pair-fiber decomposition of degree-two energies]
\label{lem:sol-frd2-pair}
Fix $i<j$ and abbreviate $s=s_{ij}$, $\delta=P_i-P_j$.  Then, conditionally on
$\mathcal F_{ij}$ and on the event $\{s>0\}$ (whose complement is $\mu$-null), the
ratio $\delta/s$ is uniform on $[-1,1]$ and independent of $\mathcal F_{ij}$;
consequently
\begin{equation}
  \Var_{ij}(\delta)=\frac{s^2}3,\qquad
  \Var_{ij}(\delta^2)=\frac{4s^4}{45},\qquad
  \Cov_{ij}(\delta,\delta^2)=0.
  \label{eq:sol-frd2-delta-moments}
\end{equation}
Every polynomial $f$ of degree at most $2$ can be written on the pair fiber, via the
substitution $p_i=\tfrac{s+\delta}2$, $p_j=\tfrac{s-\delta}2$, uniquely as
\begin{equation}
  f=\alpha_f+\beta_f\,\delta+\gamma_f\,\delta^2,
  \label{eq:sol-frd2-abc}
\end{equation}
where $\alpha_f,\beta_f,\gamma_f$ are polynomials in $s$ and $(p_\ell)_{\ell\ne i,j}$
(hence $\mathcal F_{ij}$-measurable), with $\deg\beta_f\le1$ and $\gamma_f$ constant.
For two such polynomials $f,g$,
\begin{equation}
  \E\,\frac{\Cov_{ij}(f,g)}{s^2}
  =\frac13\,\E\bigl[\beta_f\beta_g\bigr]
  +\frac4{45}\,\E\bigl[\gamma_f\gamma_g\,s^2\bigr],
  \label{eq:sol-frd2-pair-energy}
\end{equation}
a finite quantity.  In particular $K$ is finite, symmetric, bilinear, positive
semidefinite, and well defined on $V_{m,2}$.
\end{lemma}

\begin{proof}
The law $\mu$ has constant density with respect to $(m-1)$-dimensional Hausdorff
measure on $\Delta_{m-1}$.  Conditioning on $(P_\ell)_{\ell\ne i,j}$ fixes
$s=1-\sum_{\ell\ne i,j}P_\ell$ and leaves $(P_i,P_j)$ with a constant conditional
density on the segment $\{(x,y):x,y\ge0,\ x+y=s\}$; parametrized by
$P_i\in[0,s]$ this is the uniform law, so $\delta/s=2P_i/s-1$ is uniform on $[-1,1]$,
with conditional law not depending on $\mathcal F_{ij}$.  The event $\{s=0\}$ is the
face $\{P_i=P_j=0\}$, which is $\mu$-null.  With $V=\delta/s$ uniform on $[-1,1]$:
$\E V=\E V^3=0$, $\E V^2=\tfrac13$, $\E V^4=\tfrac15$, whence
$\Var(\delta\mid\mathcal F_{ij})=s^2/3$,
$\Var(\delta^2\mid\mathcal F_{ij})=s^4(\tfrac15-\tfrac19)=\tfrac{4s^4}{45}$, and
$\Cov(\delta,\delta^2\mid\mathcal F_{ij})=s^3\,\E V^3=0$, proving
\eqref{eq:sol-frd2-delta-moments}.

For \eqref{eq:sol-frd2-abc}, substitute
$p_i=\tfrac{s+\delta}2$, $p_j=\tfrac{s-\delta}2$ into each monomial of $f$:
\begin{equation}
\begin{gathered}
  p_i=\tfrac s2+\tfrac\delta2,\qquad
  p_j=\tfrac s2-\tfrac\delta2,\qquad
  p_i^2=\tfrac{s^2}4+\tfrac s2\,\delta+\tfrac14\,\delta^2,\qquad
  p_j^2=\tfrac{s^2}4-\tfrac s2\,\delta+\tfrac14\,\delta^2,\\
  p_ip_j=\tfrac{s^2}4-\tfrac14\,\delta^2,\qquad
  p_ip_\ell=\tfrac{p_\ell}2\,s+\tfrac{p_\ell}2\,\delta,\qquad
  p_jp_\ell=\tfrac{p_\ell}2\,s-\tfrac{p_\ell}2\,\delta
  \quad(\ell\ne i,j),
\end{gathered}
  \label{eq:sol-frd2-substitution-table}
\end{equation}
while monomials not involving $p_i,p_j$ contribute to $\alpha_f$ only.  Collecting
powers of $\delta$ gives \eqref{eq:sol-frd2-abc} with the stated degrees; uniqueness
holds because $(1,\delta,\delta^2)$ are linearly independent functions of $\delta$ on
any fiber with $s>0$.  Since $\alpha_f$ is $\mathcal F_{ij}$-measurable, bilinearity
of conditional covariance and \eqref{eq:sol-frd2-delta-moments} give
\[
  \Cov_{ij}(f,g)
  =\beta_f\beta_g\,\frac{s^2}3+\gamma_f\gamma_g\,\frac{4s^4}{45},
\]
the two cross terms carrying the vanishing factor
$\Cov_{ij}(\delta,\delta^2)=0$.  Dividing by $s^2$ on $\{s>0\}$ and taking
expectations proves \eqref{eq:sol-frd2-pair-energy}; the integrands are polynomials,
hence bounded on the simplex, so the expectation is finite.  Positive semidefiniteness
of $K$ is clear from $\Var_{ij}\ge0$.  Conditional covariances depend only on the
a.e.\ classes of $f,g$ and vanish when one argument is a.e.\ constant, so $K$ descends
to $V_{m,2}$.
\end{proof}

\subsection*{Symmetry structure of the degree-two quotient}

Throughout, ``module'' means a finite-dimensional real representation of $S_m$; all
maps are $S_m$-equivariant unless stated otherwise.  We use the following standard
facts about a finite group $\Gamma$ acting on real vector spaces; each has a
one-paragraph classical proof, recalled for self-containment.
\begin{itemize}
\item[(F1)] (Complete reducibility.)  Averaging any inner product over $\Gamma$
  produces an invariant inner product; the orthogonal complement of a submodule is a
  submodule; hence every module is a direct sum of irreducibles, every submodule has
  an invariant complement, and every quotient of a module is isomorphic to a
  submodule.  Consequently, in a short exact sequence of modules the multiplicity of
  each irreducible is additive.
\item[(F2)] (Schur.)  A nonzero map between irreducibles is an isomorphism; the image
  of an irreducible submodule under any equivariant map is $0$ or isomorphic to it.
  The \emph{$E$-isotypic component} $W_E(V)$ of a module $V$ is the sum of all
  submodules isomorphic to the irreducible $E$; equivariant maps send
  $E$-isotypic parts into $E$-isotypic parts, and $V$ is the direct sum of its
  isotypic components.
\item[(F3)] (Orbit counting for permutation modules.)  If $\Gamma$ acts on finite
  sets $S,T$, the space of equivariant linear maps $\R^S\to\R^T$ has dimension equal
  to the number of $\Gamma$-orbits on $T\times S$: an equivariant matrix is exactly a
  matrix constant on orbits.
\item[(F4)] (Multiplicity via Hom.)  If $E$ is irreducible with
  $\operatorname{End}_\Gamma(E)=\R\,\mathrm{id}$, then for every module $V$ the
  multiplicity of $E$ in $V$ equals $\dim\operatorname{Hom}_\Gamma(E,V)$, and
  $W_E(V)\cong E\otimes\R^{k}$ with $k$ that multiplicity.  (Decompose $V$ into
  irreducibles by (F1); by (F2) and Schur, $\operatorname{Hom}_\Gamma(E,\cdot)$ of each
  summand is $\R$ for summands $\cong E$ and $0$ otherwise.)  In particular the
  multiplicity of $E$ is monotone under submodules and quotients.
\item[(F5)] (Invariant bilinear forms on an irreducible.)  If $E$ is irreducible with
  $\operatorname{End}_\Gamma(E)=\R\,\mathrm{id}$, the space of invariant bilinear
  forms on $E$ is one-dimensional, spanned by any fixed invariant inner product
  $J_E$: a form corresponds to an equivariant map $E\to E^*\cong E$, i.e.\ to an
  element of $\operatorname{End}_\Gamma(E)=\R$.
\end{itemize}

\begin{lemma}[the three irreducibles]
\label{lem:sol-frd2-irreducibles}
Let $M=\R^m$ with permuted coordinates and, for $m\ge3$, let $N=\R^{\binom{[m]}2}$ be
the permutation module on unordered pairs $\{i,j\}$, with basis $(x_{ij})_{i<j}$.
Then:
\begin{enumerate}
\item $M=\mathrm{triv}\oplus\mathrm{std}$, where $\mathrm{triv}=\R\one$ and
  $\mathrm{std}=\{a:\sum_ia_i=0\}$ is irreducible of dimension $m-1$ with
  $\operatorname{End}_{S_m}(\mathrm{std})=\R\,\mathrm{id}$, and
  $\mathrm{std}\not\cong\mathrm{triv}$.
\item For $m\ge4$, $N=\mathrm{triv}\oplus\mathrm{std}\oplus\mathrm X$ with each
  summand of multiplicity one, where $\mathrm X$ is irreducible of dimension
  $m(m-3)/2$, satisfies $\operatorname{End}_{S_m}(\mathrm X)=\R\,\mathrm{id}$, and
  $\mathrm X\not\cong\mathrm{triv},\mathrm{std}$.  For $m=3$,
  $N=\mathrm{triv}\oplus\mathrm{std}$ and the summand $\mathrm X$ is absent.
\end{enumerate}
\end{lemma}

\begin{proof}
(1)  $S_m$ has two orbits on $[m]^2$ (diagonal, off-diagonal), so
$\dim\operatorname{End}(M)=2$ by (F3).  Also $M^{S_m}=\R\one$ is one-dimensional, so
$\mathrm{triv}$ has multiplicity one; let $C=\one^\perp$ be its invariant complement.
Then $C^{S_m}=0$, so $\operatorname{Hom}(\mathrm{triv},C)=0$ and
$2=\dim\operatorname{End}(M)
 =\dim\operatorname{End}(\mathrm{triv})+\dim\operatorname{End}(C)
 =1+\dim\operatorname{End}(C)$.
A decomposable module has endomorphism algebra of dimension at least $2$ (the two
projections), so $\operatorname{End}(C)=\R\,\mathrm{id}$ forces $C$ irreducible.
Finally $C\not\cong\mathrm{triv}$ because $C^{S_m}=0$.  Set $\mathrm{std}=C$.

(2)  Let $m\ge4$.  The orbits of $S_m$ on (ordered) pairs of $2$-subsets are
classified by $\abs{A\cap B}\in\{2,1,0\}$, all realized, so
$\dim\operatorname{End}(N)=3$.  The orbits on $[m]\times\binom{[m]}2$ are classified
by $i\in A$ versus $i\notin A$, so $\dim\operatorname{Hom}(M,N)=2$.  As
$N^{S_m}$ is one-dimensional (one orbit of $2$-subsets), $\mathrm{triv}$ has
multiplicity one in $N$ and
$\dim\operatorname{Hom}(\mathrm{std},N)
 =\dim\operatorname{Hom}(M,N)-\dim\operatorname{Hom}(\mathrm{triv},N)=2-1=1$,
so by (F4) $\mathrm{std}$ has multiplicity exactly one.  Decompose
$N=\mathrm{triv}\oplus S'\oplus \mathrm X$ with $S'\cong\mathrm{std}$ and $\mathrm X$
the sum of the remaining irreducible summands, which contains no copy of
$\mathrm{triv}$ or $\mathrm{std}$.  Then cross-Homs between the three summands vanish
and
$3=\dim\operatorname{End}(N)=1+1+\dim\operatorname{End}(\mathrm X)$,
so as in (1) $\mathrm X$ is irreducible with
$\operatorname{End}(\mathrm X)=\R\,\mathrm{id}$, and
$\mathrm X\not\cong\mathrm{triv},\mathrm{std}$ (otherwise those multiplicities would
exceed one).  Its dimension is $\binom m2-1-(m-1)=m(m-3)/2$.  For $m=3$ there is no
orbit with $\abs{A\cap B}=0$, so $\dim\operatorname{End}(N)=2$ and the same argument
yields $N=\mathrm{triv}\oplus\mathrm{std}$ (dimensions: $3=1+2$).
\end{proof}

\begin{lemma}[isotypic decomposition of $V_{m,2}$]
\label{lem:sol-frd2-decomposition}
For $m\ge3$, as $S_m$-modules,
\begin{equation}
  V_{m,2}\;\cong\;
  \mathrm{triv}\ \oplus\ \mathrm{std}\otimes\R^2\ \oplus\ \mathrm X,
  \qquad
  \dim V_{m,2}=\frac{(m-1)(m+2)}2,
  \label{eq:sol-frd2-isotypic}
\end{equation}
where the $\mathrm X$-summand is absent when $m=3$.  Moreover:
\begin{enumerate}
\item the $\mathrm{triv}$-isotypic component is
  $W_{\rm triv}=\R\,[f_0]$ with $f_0=\sum_kp_k^2$, and $[f_0]\ne0$;
\item the classes $L=[p_1-p_2]$ and $Q=[p_1^2-p_2^2]$ are nonzero, linearly
  independent, and lie in the $\mathrm{std}$-isotypic component
  $W_{\rm std}$; they are the images of the common seed vector
  $v_0=e_1-e_2\in\mathrm{std}$ under the two equivariant maps
  $T_1:a\mapsto[\textstyle\sum_ia_ip_i]$ and
  $T_2:a\mapsto[\textstyle\sum_ia_ip_i^2]$, whose images span $W_{\rm std}$;
\item for $m\ge4$, the class $F=[(p_1-p_3)(p_2-p_4)]$ is nonzero and lies in the
  $\mathrm X$-isotypic component $W_{\mathrm X}$, which it generates.
\end{enumerate}
\end{lemma}

\begin{proof}
\emph{Kernel of the quotient map.}
Define the equivariant surjection $\Psi:W_{\rm poly}\to V_{m,2}$, $f\mapsto[f]$.  We
claim
\begin{equation}
  \ker\Psi
  =\Bigl\{\,b_0\textstyle\sum_ip_i
    +\bigl(\textstyle\sum_ip_i-1\bigr)\textstyle\sum_ib_ip_i
    \;:\;b_0\in\R,\ b\in\R^m\Bigr\}
  \;\cong\;\mathrm{triv}\oplus M
  \;\cong\;2\,\mathrm{triv}\oplus\mathrm{std}.
  \label{eq:sol-frd2-kernel}
\end{equation}
Indeed $f\in\ker\Psi$ means $f=c$ $\mu$-a.e.\ for some $c\in\R$.  The support
$\Delta_{m-1}$ has nonempty relative interior in the hyperplane
$H=\{\sum_ip_i=1\}$, and a polynomial vanishing on a relatively open subset of an
affine subspace vanishes on the whole subspace; hence $f-c$ vanishes on $H$.
Choosing affine coordinates in which $t=\sum_ip_i-1$ is one coordinate and dividing
by $t$ (polynomial division in $t$, degree bookkeeping: $\deg\le2$), we get
$f=c+(\sum_ip_i-1)\ell$ with $\ell$ affine, say $\ell=b_0+\sum_ib_ip_i$.  Since $f$
has zero constant term, $c=b_0$, which yields exactly the displayed set; conversely
every element of that set is $\mu$-a.e.\ constant (equal to $b_0$).  The
parametrization $(b_0,b)\mapsto f$ is linear, injective (compare first the quadratic
homogeneous part $(\sum_ip_i)(\sum_ib_ip_i)$, which vanishes only if $b=0$ since the
polynomial ring is a domain; then $b_0\sum p_i=0$ forces $b_0=0$), and equivariant
for the action fixing $b_0$ and permuting $b$.  This proves
\eqref{eq:sol-frd2-kernel}.

\emph{Multiplicities.}
The monomial families $(p_i)$, $(p_i^2)$, $(p_ip_j)_{i<j}$ are linearly independent
in the polynomial ring and are permuted by $S_m$ exactly as the bases of $M$, $M$,
$N$; hence $W_{\rm poly}\cong M\oplus M\oplus N$, whose isotypic multiplicities are
(by Lemma~\ref{lem:sol-frd2-irreducibles}): $\mathrm{triv}$: $3$, $\mathrm{std}$:
$3$, $\mathrm X$: $1$ (for $m\ge4$; $\mathrm X$ absent at $m=3$).  By (F1) applied to
$0\to\ker\Psi\to W_{\rm poly}\to V_{m,2}\to0$ and \eqref{eq:sol-frd2-kernel}, the
multiplicities in $V_{m,2}$ are $\mathrm{triv}$: $3-2=1$, $\mathrm{std}$: $3-1=2$,
$\mathrm X$: $1-0=1$, and no other irreducible occurs.  This is
\eqref{eq:sol-frd2-isotypic}; the dimension is
$1+2(m-1)+m(m-3)/2=(m-1)(m+2)/2$ (also valid at $m=3$).

\emph{(1).}
$[f_0]$ is $S_m$-fixed.  It is nonzero: if $f_0\in\ker\Psi$, comparing quadratic
homogeneous parts in \eqref{eq:sol-frd2-kernel} gives
$\sum_ip_i^2=(\sum_ip_i)(\sum_ib_ip_i)$; the coefficient of $p_i^2$ forces $b_i=1$
for all $i$, but then the coefficient of $p_1p_2$ on the right is $2\ne0$, a
contradiction.  Since the $\mathrm{triv}$-multiplicity is one,
$W_{\rm triv}=\R[f_0]$.

\emph{(2).}
$T_1,T_2$ are restrictions to $\mathrm{std}\subset M$ of the equivariant maps
$a\mapsto[\sum a_ip_i]$, $a\mapsto[\sum a_ip_i^2]$, so their images lie in
$W_{\rm std}$ by (F2), and $L=T_1v_0$, $Q=T_2v_0$.  Linear independence of $L,Q$ in
$V_{m,2}$: suppose $aL+bQ\in\ker\Psi$, i.e.
$a(p_1-p_2)+b(p_1^2-p_2^2)
 =b_0\sum_ip_i+(\sum_ip_i-1)\sum_ib_ip_i$.
Quadratic parts: $b(p_1^2-p_2^2)=(\sum_ip_i)(\sum_ib_ip_i)$; the $p_i^2$
coefficients give $b_1=b$, $b_2=-b$, $b_\ell=0$ ($\ell\ge3$), and then the
$p_1p_3$ coefficient (using $m\ge3$) gives $0=b_1+b_3=b$, so $b=0$ and $b\equiv0$.
The linear part then reads $a(p_1-p_2)=b_0\sum_ip_i$, forcing $a=b_0=-a$, i.e.\
$a=0$.  In particular $L,Q\ne0$, so $T_1,T_2\ne0$ and, $\mathrm{std}$ being
irreducible, $T_1(\mathrm{std})\cong T_2(\mathrm{std})\cong\mathrm{std}$.  The two
image submodules are distinct: if $T_1(\mathrm{std})=T_2(\mathrm{std})=:U$, then
$T_1,T_2\in\operatorname{Hom}(\mathrm{std},U)$, a one-dimensional space by (F4)--(F5)
(as $U\cong\mathrm{std}$), so $T_2=cT_1$ and $Q=cL$, contradicting independence.
Distinct irreducible submodules intersect in $0$, so
$T_1(\mathrm{std})\oplus T_2(\mathrm{std})\subseteq W_{\rm std}$ has dimension
$2(m-1)=\dim W_{\rm std}$; hence the images span $W_{\rm std}$.

\emph{(3).}
Let $m\ge4$ and let $\psi:N\to V_{m,2}$ be the equivariant map
$x_{ij}\mapsto[p_ip_j]$.  Consider
\[
  v=x_{12}-x_{23}-x_{14}+x_{34}\in N,
  \qquad
  \psi(v)=[p_1p_2-p_2p_3-p_1p_4+p_3p_4]=[(p_1-p_3)(p_2-p_4)]=F.
\]
We check $v\in W_{\mathrm X}(N)$.  With the standard (invariant) inner product on
$N$, the isotypic components of $N$ are mutually orthogonal.  The
$\mathrm{triv}$-component is spanned by $\sum_{i<j}x_{ij}$, and the
$\mathrm{std}$-component lies in the image of the equivariant map $A:M\to N$,
$e_k\mapsto u_k:=\sum_{\ell\ne k}x_{k\ell}$ (indeed $A(\mathrm{std})\ne0$ because
$u_1-u_2$ has coefficient $1$ on $x_{13}$, and the $\mathrm{std}$-multiplicity of $N$
is one, so $W_{\rm std}(N)=A(\mathrm{std})$).  Now
$\langle v,\sum_{i<j}x_{ij}\rangle=1-1-1+1=0$ and, for every $k$,
$\langle v,u_k\rangle=0$: for $k=1$ the pairs of $v$ containing $1$ are
$x_{12}$ ($+$) and $x_{14}$ ($-$); for $k=2$: $x_{12}$ ($+$), $x_{23}$ ($-$); for
$k=3$: $x_{23}$ ($-$), $x_{34}$ ($+$); for $k=4$: $x_{14}$ ($-$), $x_{34}$ ($+$);
for $k\ge5$: none.  Hence $v\perp W_{\rm triv}(N)\oplus W_{\rm std}(N)$, i.e.\
$v\in W_{\mathrm X}(N)$.  By (F2), $F=\psi(v)\in W_{\mathrm X}(V_{m,2})$.  Finally
$F\ne0$: if $(p_1-p_3)(p_2-p_4)\in\ker\Psi$, comparing quadratic parts with
\eqref{eq:sol-frd2-kernel} gives
$(p_1-p_3)(p_2-p_4)=(\sum_ip_i)(\sum_ib_ip_i)$; the left side has no $p_i^2$ terms,
so all $b_i=0$, but the left side has $p_1p_2$-coefficient $1$, a contradiction.
Since $W_{\mathrm X}$ has multiplicity one and $\psi|_{W_{\mathrm X}(N)}$ is
injective (its restriction to the irreducible $W_{\mathrm X}(N)$ is nonzero), $F$
generates $W_{\mathrm X}$ as a module.
\end{proof}

\begin{lemma}[invariant pencil reduction]
\label{lem:sol-frd2-reduction}
Let $V$ be a module with invariant symmetric bilinear forms $B$ (arbitrary) and $G$
(positive definite), and let $V=\bigoplus_EW_E$ be its isotypic decomposition over
pairwise non-isomorphic irreducibles $E$ with
$\operatorname{End}(E)=\R\,\mathrm{id}$.  Then:
\begin{enumerate}
\item $B(W_E,W_{E'})=0$ for $E\not\cong E'$; hence $B$ is positive semidefinite
  (resp.\ definite) on $V$ if and only if it is so on each $W_E$.
\item Suppose $W_E$ is spanned by the images of equivariant maps
  $T_1,\dots,T_k:E\to V$ that are linearly independent in
  $\operatorname{Hom}_{S_m}(E,V)$, and fix $0\ne v_0\in E$.  Then the map
  $\Phi:E\otimes\R^k\to W_E$, $u\otimes\varepsilon_a\mapsto T_au$, is an
  isomorphism, and there is a symmetric $k\times k$ matrix $\widehat B$ with
  \begin{equation}
    B(T_au,T_bu')=\widehat B_{ab}\,J_E(u,u')
    \qquad(u,u'\in E,\ 1\le a,b\le k),
    \label{eq:sol-frd2-tensor-form}
  \end{equation}
  where $J_E$ is a fixed invariant inner product on $E$.  Consequently $B$ is
  positive semidefinite (resp.\ definite) on $W_E$ if and only if the $k\times k$
  Gram matrix $\bigl(B(T_av_0,T_bv_0)\bigr)_{a,b}=J_E(v_0,v_0)\,\widehat B$ is
  positive semidefinite (resp.\ definite).
\end{enumerate}
\end{lemma}

\begin{proof}
(1)  The invariant form $B$ and an invariant inner product on $W_{E'}$ induce an
equivariant map $W_E\to W_{E'}^*\cong W_{E'}$.  By (F2) and Schur, every equivariant
map from an $E$-isotypic to an $E'$-isotypic module is zero when
$E\not\cong E'$.  Hence $B(W_E,W_{E'})=0$, and $B$ decomposes as the orthogonal
(with respect to the decomposition) sum of its restrictions.

(2)  $\Phi$ is equivariant and surjective (images span), and
$\dim(E\otimes\R^k)=k\dim E=\dim W_E$ because the multiplicity of $E$ in $W_E$
equals $\dim\operatorname{Hom}(E,W_E)\ge k$ by linear independence, while
surjectivity gives $\le k$; so $\Phi$ is an isomorphism.  For fixed $a,b$, the
bilinear form $(u,u')\mapsto B(T_au,T_bu')$ on $E$ is invariant, hence by (F5) equals
$\widehat B_{ab}J_E(u,u')$ for a unique scalar $\widehat B_{ab}$; symmetry of $B$ and
of $J_E$ gives $\widehat B_{ab}=\widehat B_{ba}$.  Now for
$f=\Phi\bigl(\sum_\alpha u^\alpha\otimes s_\alpha\bigr)$ with $(u^\alpha)_\alpha$ a
$J_E$-orthonormal basis of $E$ and $s_\alpha\in\R^k$,
\begin{equation}
  B(f,f)
  =\sum_{\alpha,\gamma}\sum_{a,b}
   s_\alpha^a s_\gamma^b\,B(T_au^\alpha,T_bu^\gamma)
  =\sum_{\alpha,\gamma}J_E(u^\alpha,u^\gamma)\,s_\alpha^T\widehat B\,s_\gamma
  =\sum_\alpha s_\alpha^T\widehat B\,s_\alpha.
  \label{eq:sol-frd2-tensor-sum}
\end{equation}
Since every element of $W_E$ is such an $f$ and the vectors
$(s_\alpha)_\alpha$ are arbitrary, $B\ge0$ (resp.\ $>0$) on $W_E$ iff
$\widehat B\succeq0$ (resp.\ $\succ0$), iff the Gram matrix
$J_E(v_0,v_0)\widehat B$ is so ($J_E(v_0,v_0)>0$).
\end{proof}

\subsection*{Sector computations}

Throughout this subsection $E_{ij}(f,g)$ denotes the pair energy
\eqref{eq:sol-frd2-pair-energy} and we tabulate the coefficients
$(\beta,\gamma)$ of Lemma~\ref{lem:sol-frd2-pair} for each seed and each pair type;
pairs not listed have $\beta=\gamma=0$.  Recall
\eqref{eq:sol-frd2-small-moments} and $D_2=m(m+1)$, $D_3=D_2(m+2)$,
$D_4=D_3(m+3)$.

\begin{lemma}[trivial sector]
\label{lem:sol-frd2-trivial}
With $f_0=\sum_kp_k^2$,
\begin{equation}
  K(f_0,f_0)=\frac{4(m-1)}{5m^2(m+1)^2},
  \qquad
  G(f_0,f_0)=\frac{4(m-1)}{(m+1)^2(m+2)(m+3)},
  \qquad
  \frac{K(f_0,f_0)}{G(f_0,f_0)}=\frac{(m+2)(m+3)}{5m^2}.
  \label{eq:sol-frd2-trivial}
\end{equation}
Moreover $[f_0]=\tfrac1{m(m+1)}\,[\abs X^2-(m-1)]$ in $V_{m,2}$.
\end{lemma}

\begin{proof}
For every pair $(i,j)$, by \eqref{eq:sol-frd2-substitution-table},
$p_i^2+p_j^2=\tfrac{s^2}2+\tfrac{\delta^2}2$, so
$(\beta,\gamma)=(0,\tfrac12)$ and, by \eqref{eq:sol-frd2-pair-energy} and
\eqref{eq:sol-frd2-small-moments},
$E_{ij}(f_0,f_0)=\tfrac4{45}\cdot\tfrac14\,\E s^2=\tfrac1{45}\cdot\tfrac6{D_2}
=\tfrac2{15D_2}$.
Summing over the $\binom m2$ pairs and inserting the prefactor of
\eqref{eq:sol-frd2-forms},
\[
  K(f_0,f_0)
  =\frac{12}{m^2(m+1)}\cdot\frac{m(m-1)}2\cdot\frac2{15\,m(m+1)}
  =\frac{4(m-1)}{5m^2(m+1)^2}.
\]
For the variance, Lemma~\ref{lem:sol-frd2-moments} gives
$\E f_0=m\cdot\tfrac2{D_2}=\tfrac2{m+1}$ and
$\E f_0^2=m\,\E P_1^4+m(m-1)\,\E P_1^2P_2^2
=\tfrac{24m+4m(m-1)}{D_4}=\tfrac{4m(m+5)}{D_4}$, so
\[
  G(f_0,f_0)
  =\frac{4(m+5)}{(m+1)(m+2)(m+3)}-\frac4{(m+1)^2}
  =\frac{4\bigl[(m+5)(m+1)-(m+2)(m+3)\bigr]}{(m+1)^2(m+2)(m+3)}
  =\frac{4(m-1)}{(m+1)^2(m+2)(m+3)},
\]
using $(m+5)(m+1)-(m+2)(m+3)=m-1$.  The quotient is
$(m+2)(m+3)/(5m^2)$.  Finally
$\abs X^2=m(m+1)\sum_i(p_i-\tfrac1m)^2=m(m+1)\bigl(f_0-\tfrac1m\bigr)$, so
$[\abs X^2-(m-1)]=m(m+1)[f_0]$.
\end{proof}

\begin{lemma}[standard sector Gram matrices]
\label{lem:sol-frd2-standard}
With $L=p_1-p_2$ and $Q=p_1^2-p_2^2$,
\begin{equation}
\begin{aligned}
  K(L,L)&=\frac2{m(m+1)}, &
  K(L,Q)&=\frac4{m^2(m+1)}, &
  K(Q,Q)&=\frac{8(8m-1)}{5m^3(m+1)^2},\\
  G(L,L)&=\frac2{m(m+1)}, &
  G(L,Q)&=\frac8{m(m+1)(m+2)}, &
  G(Q,Q)&=\frac{40}{m(m+1)(m+2)(m+3)}.
\end{aligned}
  \label{eq:sol-frd2-standard-gram}
\end{equation}
\end{lemma}

\begin{proof}
\emph{Energies.}  The $(\beta,\gamma)$ tables from
\eqref{eq:sol-frd2-substitution-table}, with $j$ always denoting an index
$\ge3$ distinct from $1,2$:
\[
\begin{array}{l|cc|cc}
 \text{pair} & \beta_L & \gamma_L & \beta_Q & \gamma_Q\\\hline
 (1,2) & 1 & 0 & s & 0\\
 (1,j) & \tfrac12 & 0 & \tfrac s2 & \tfrac14\\
 (2,j) & -\tfrac12 & 0 & -\tfrac s2 & -\tfrac14
\end{array}
\]
(For $(1,j)$: $L=\tfrac{s+\delta}2-p_2$ and
$Q=\tfrac{s^2}4+\tfrac s2\delta+\tfrac{\delta^2}4-p_2^2$; for $(2,j)$, with
$\delta=p_2-p_j$: $L=p_1-\tfrac{s+\delta}2$,
$Q=p_1^2-\tfrac{s^2}4-\tfrac s2\delta-\tfrac{\delta^2}4$; for $(1,2)$:
$L=\delta$, $Q=s\delta$.)  By
\eqref{eq:sol-frd2-pair-energy}--\eqref{eq:sol-frd2-small-moments}:
\begin{align*}
  E_{12}(L,L)&=\tfrac13, &
  E_{1j}(L,L)=E_{2j}(L,L)&=\tfrac1{12},\\
  E_{12}(L,Q)&=\tfrac13\E s=\tfrac2{3m}, &
  E_{1j}(L,Q)=E_{2j}(L,Q)&=\tfrac13\cdot\tfrac{\E s}4=\tfrac1{6m},\\
  E_{12}(Q,Q)&=\tfrac13\E s^2=\tfrac2{D_2}, &
  E_{1j}(Q,Q)=E_{2j}(Q,Q)
  &=\tfrac13\cdot\tfrac{\E s^2}4+\tfrac4{45}\cdot\tfrac{\E s^2}{16}
   =\tfrac{16}{180}\,\E s^2=\tfrac8{15D_2}.
\end{align*}
There are $2(m-2)$ pairs of the types $(1,j),(2,j)$, so
\begin{align*}
  \sum_{i<j}E_{ij}(L,L)&=\tfrac13+\tfrac{2(m-2)}{12}=\tfrac m6,\\
  \sum_{i<j}E_{ij}(L,Q)&=\tfrac2{3m}+\tfrac{2(m-2)}{6m}=\tfrac{2+(m-2)}{3m}
   =\tfrac13,\\
  \sum_{i<j}E_{ij}(Q,Q)&=\tfrac2{D_2}+\tfrac{16(m-2)}{15D_2}
   =\tfrac{2(8m-1)}{15D_2}.
\end{align*}
Multiplying by $12/(m^2(m+1))$ gives the three $K$-entries in
\eqref{eq:sol-frd2-standard-gram}.

\emph{Variances.}  By symmetry $\E L=\E Q=0$, so, by
Lemma~\ref{lem:sol-frd2-moments},
\begin{align*}
  G(L,L)&=\E(P_1-P_2)^2=\tfrac2{D_2},\\
  G(L,Q)&=\E\bigl[(P_1-P_2)(P_1^2-P_2^2)\bigr]
   =\E\bigl[P_1^3+P_2^3-P_1P_2(P_1+P_2)\bigr]
   =\tfrac{12-4}{D_3}=\tfrac8{D_3},\\
  G(Q,Q)&=\E(P_1^2-P_2^2)^2
   =2\,\E P_1^4-2\,\E P_1^2P_2^2=\tfrac{48-8}{D_4}=\tfrac{40}{D_4}.
  \qedhere
\end{align*}
\end{proof}

\begin{lemma}[two-row sector, $m\ge4$]
\label{lem:sol-frd2-tworow}
With $F=(p_1-p_3)(p_2-p_4)$,
\begin{equation}
  K(F,F)=\frac{8(5m-4)}{5m^3(m+1)^2},
  \qquad
  G(F,F)=\frac4{m(m+1)(m+2)(m+3)},
  \qquad
  \frac{K(F,F)}{G(F,F)}
  =\frac{2(5m-4)(m+2)(m+3)}{5m^2(m+1)}.
  \label{eq:sol-frd2-tworow}
\end{equation}
\end{lemma}

\begin{proof}
\emph{Energy.}  The $(\beta,\gamma)$ table, computed from
\eqref{eq:sol-frd2-substitution-table} (here $j\ge5$; in each row $\delta=p_i-p_j$
for the listed pair $(i,j)$, $i<j$):
\[
\begin{array}{l|cc|l}
 \text{pair} & \beta_F & \gamma_F & \text{derivation}\\\hline
 (1,3) & p_2-p_4 & 0 & F=\delta\,(p_2-p_4)\\
 (2,4) & p_1-p_3 & 0 & F=(p_1-p_3)\,\delta\\
 (1,2) & \tfrac{p_3-p_4}2 & -\tfrac14 &
  F=\bigl(\tfrac{s+\delta}2-p_3\bigr)\bigl(\tfrac{s-\delta}2-p_4\bigr)\\
 (3,4) & \tfrac{p_1-p_2}2 & -\tfrac14 &
  F=\bigl(p_1-\tfrac{s+\delta}2\bigr)\bigl(p_2-\tfrac{s-\delta}2\bigr)\\
 (1,4) & \tfrac{p_2-p_3}2 & \tfrac14 &
  F=\bigl(\tfrac{s+\delta}2-p_3\bigr)\bigl(p_2-\tfrac{s-\delta}2\bigr)\\
 (2,3) & \tfrac{p_1-p_4}2 & \tfrac14 &
  F=\bigl(p_1-\tfrac{s-\delta}2\bigr)\bigl(\tfrac{s+\delta}2-p_4\bigr)\\
 (1,j) & \tfrac{p_2-p_4}2 & 0 & F=\bigl(\tfrac{s+\delta}2-p_3\bigr)(p_2-p_4)\\
 (2,j) & \tfrac{p_1-p_3}2 & 0 & \\
 (3,j) & -\tfrac{p_2-p_4}2 & 0 & \\
 (4,j) & -\tfrac{p_1-p_3}2 & 0 &
\end{array}
\]
For instance, for the pair $(1,2)$:
$F=\bigl[(\tfrac s2-p_3)+\tfrac\delta2\bigr]
   \bigl[(\tfrac s2-p_4)-\tfrac\delta2\bigr]
 =\alpha+\tfrac{p_3-p_4}2\,\delta-\tfrac14\,\delta^2$
with $\alpha$ $\mathcal F_{12}$-measurable; the rows $(3,4)$, $(1,4)$, $(2,3)$ are
identical expansions.  By
\eqref{eq:sol-frd2-pair-energy}--\eqref{eq:sol-frd2-small-moments},
\begin{align*}
  E_{13}(F,F)=E_{24}(F,F)
  &=\tfrac13\,\E(P_2-P_4)^2=\tfrac2{3D_2},\\
  E_{12}(F,F)=E_{34}(F,F)=E_{14}(F,F)=E_{23}(F,F)
  &=\tfrac13\cdot\tfrac{\E(P_3-P_4)^2}4+\tfrac4{45}\cdot\tfrac{\E s^2}{16}
   =\tfrac1{6D_2}+\tfrac1{30D_2}=\tfrac1{5D_2},\\
  E_{ij}(F,F)\ \ (i\in\{1,2,3,4\},\ j\ge5)
  &=\tfrac13\cdot\tfrac{\E(P_a-P_b)^2}4=\tfrac1{6D_2}
  \qquad(4(m-4)\text{ pairs}).
\end{align*}
Hence
\[
  \sum_{i<j}E_{ij}(F,F)
  =\frac1{D_2}\Bigl[\frac43+\frac45+\frac{4(m-4)}6\Bigr]
  =\frac1{D_2}\cdot\frac{20+12+10(m-4)}{15}
  =\frac{2(5m-4)}{15\,D_2},
\]
and $K(F,F)=\tfrac{12}{m^2(m+1)}\cdot\tfrac{2(5m-4)}{15m(m+1)}
=\tfrac{8(5m-4)}{5m^3(m+1)^2}$.

\emph{Variance.}  $\E F=\E P_1P_2-\E P_1P_4-\E P_2P_3+\E P_3P_4=0$, and expanding
$F^2=(p_1^2-2p_1p_3+p_3^2)(p_2^2-2p_2p_4+p_4^2)$ termwise with
Lemma~\ref{lem:sol-frd2-moments} (all four indices distinct),
\[
  \E F^2
  =\frac{4\cdot4-8\cdot2+4\cdot1}{D_4}
  =\frac4{D_4},
\]
the three groups being: the four square--square products, each with
$\E P_a^2P_b^2=4/D_4$, totalling $16/D_4$; the four products of a square with a
$-2p_bp_c$ factor, each contributing $-2\,\E P_a^2P_bP_c=-4/D_4$, totalling
$-16/D_4$; and the single product
$(-2p_1p_3)(-2p_2p_4)$, contributing $4\,\E P_1P_2P_3P_4=4/D_4$.  The quotient follows by
$G(F,F)=\E F^2$ and $D_4=m(m+1)(m+2)(m+3)$.
\end{proof}

\subsection*{Proof of Theorem~\ref{thm:sol-fiber-root-degree-two}}

\begin{proof}
Write $\lambda^*=(m+2)(m+3)/(5m^2)$ and let $\Phi=K-\lambda^*G$, an
$S_m$-invariant symmetric bilinear form on $V_{m,2}$
(Lemma~\ref{lem:sol-frd2-pair}).  By
Lemma~\ref{lem:sol-frd2-decomposition} and
Lemma~\ref{lem:sol-frd2-reduction}(1),
\[
  \Phi(f,f)=\Phi(f_{\rm triv},f_{\rm triv})
   +\Phi(f_{\rm std},f_{\rm std})+\Phi(f_{\mathrm X},f_{\mathrm X})
\]
for the isotypic components of any $f\in V_{m,2}$ (the $\mathrm X$-term absent at
$m=3$).  We treat the three sectors.

\emph{Trivial sector.}  $W_{\rm triv}=\R[f_0]$ and, by
Lemma~\ref{lem:sol-frd2-trivial},
$K(f_0,f_0)=\lambda^*G(f_0,f_0)$, i.e.\ $\Phi\equiv0$ on $W_{\rm triv}$, with
$G(f_0,f_0)>0$.

\emph{Standard sector.}  Apply Lemma~\ref{lem:sol-frd2-reduction}(2) with
$E=\mathrm{std}$, $k=2$, the maps $T_1,T_2$ and seed $v_0=e_1-e_2$ of
Lemma~\ref{lem:sol-frd2-decomposition}(2) (linearly independent in
$\operatorname{Hom}$, since $c_1T_1+c_2T_2=0$ evaluated at $v_0$ gives
$c_1L+c_2Q=0$, hence $c_1=c_2=0$).  Positive definiteness of $\Phi$ on
$W_{\rm std}$ is therefore equivalent to positive definiteness of the $2\times2$
Gram matrix of $(L,Q)$ under $\Phi$.  Multiplying the entries of
\eqref{eq:sol-frd2-standard-gram} by the common positive factor
$m(m+1)/2$ yields the normalized matrices
\[
  \widetilde K=
  \begin{pmatrix}
    1 & \tfrac2m\\[2pt]
    \tfrac2m & \tfrac{4(8m-1)}{5m^2(m+1)}
  \end{pmatrix},
  \qquad
  \widetilde G=
  \begin{pmatrix}
    1 & \tfrac4{m+2}\\[2pt]
    \tfrac4{m+2} & \tfrac{20}{(m+2)(m+3)}
  \end{pmatrix},
\]
and a direct computation gives
\begin{align*}
  (\widetilde K-\lambda^*\widetilde G)_{11}
  &=1-\frac{(m+2)(m+3)}{5m^2}
   =\frac{4m^2-5m-6}{5m^2}=\frac{(4m+3)(m-2)}{5m^2},\\
  (\widetilde K-\lambda^*\widetilde G)_{12}
  &=\frac2m-\frac{4(m+3)}{5m^2}
   =\frac{10m-4m-12}{5m^2}=\frac{6(m-2)}{5m^2},\\
  (\widetilde K-\lambda^*\widetilde G)_{22}
  &=\frac{4(8m-1)}{5m^2(m+1)}-\frac4{m^2}
   =\frac{4(8m-1)-20(m+1)}{5m^2(m+1)}=\frac{12(m-2)}{5m^2(m+1)},
\end{align*}
that is,
\begin{equation}
  5m^2\,(\widetilde K-\lambda^*\widetilde G)
  =(m-2)
  \begin{pmatrix}
    4m+3 & 6\\[2pt]
    6 & \tfrac{12}{m+1}
  \end{pmatrix}.
  \label{eq:sol-frd2-std-psd}
\end{equation}
For $m\ge3$ the factor $m-2$ is positive, the diagonal entries are positive, and
the determinant of the bracketed matrix is
$\tfrac{12(4m+3)}{m+1}-36=\tfrac{12(4m+3)-36(m+1)}{m+1}=\tfrac{12m}{m+1}>0$.
Hence $\Phi$ is positive definite on $W_{\rm std}$:
$\Phi(f_{\rm std},f_{\rm std})>0$ whenever $f_{\rm std}\ne0$.

\emph{Two-row sector ($m\ge4$).}  $W_{\mathrm X}$ has multiplicity one and is
generated by $F$ (Lemma~\ref{lem:sol-frd2-decomposition}(3)), so
Lemma~\ref{lem:sol-frd2-reduction}(2) with $k=1$ reduces positivity on
$W_{\mathrm X}$ to the sign of the scalar $\Phi(F,F)$.  By
Lemma~\ref{lem:sol-frd2-tworow},
\[
  \frac{K(F,F)}{G(F,F)}
  =\frac{2(5m-4)}{m+1}\,\lambda^*
  =\lambda^*+\frac{9(m-1)}{m+1}\,\lambda^*
  >\lambda^*,
\]
since $2(5m-4)-(m+1)=9(m-1)>0$.  Hence
$\Phi(F,F)=G(F,F)\bigl(\tfrac{K(F,F)}{G(F,F)}-\lambda^*\bigr)>0$ and $\Phi$ is
positive definite on $W_{\mathrm X}$.

\emph{Assembly.}  For every $f\in V_{m,2}$,
$\Phi(f,f)\ge0$, i.e.\ $K(f,f)\ge\lambda^*G(f,f)$, with equality if and only if
$f_{\rm std}=0$ and $f_{\mathrm X}=0$, i.e.\ $f\in W_{\rm triv}=\R[f_0]$.  Since
$K(f_0,f_0)=\lambda^*G(f_0,f_0)$ with $G(f_0,f_0)>0$, the minimum
\eqref{eq:sol-frd2-lmin} equals $\lambda^*$ and is attained exactly on
$\R[f_0]=\R[\abs X^2-(m-1)]$ (Lemma~\ref{lem:sol-frd2-trivial}).  Finally
$(m+2)(m+3)>m^2$ gives $\lambda^*>\tfrac15$ for all $m$, and
$\lambda^*\to\tfrac15$.
\end{proof}

\subsection*{Proof of Corollary~\ref{cor:sol-frd2-no-degree-two-refuter}}

\begin{proof}
All statements are conditional on the identification
\eqref{eq:sol-frd2-root-identification}, certified in
\texttt{solutions/conditional-fiber-frame-structure.tex}: the root frame
$\rho_{\rm root}$ is an admissible even frame, and for every $f$ in the maximal
form domain (in particular for every polynomial) the root conditional-fiber form
equals the pair-redistribution expression \eqref{eq:sol-frd2-forms}.

\emph{(1).}  Let $(c_r,w_r,M,\eps)$ be a degree-$2$ dual certificate
\eqref{eq:sol-frd2-certificate}.  Apply the direction bound in
\eqref{eq:sol-frd2-certificate} at the $m(m-1)$ atoms of $\rho_{\rm root}$ and
average:
\[
  \sum_rw_r\,K(c_r,c_r)
  =(m-1)\int\sum_rw_r\,\qJac_\theta[c_r]\dd\rho_{\rm root}(\theta)
  \le(m-1)\int\theta^TM\theta\dd\rho_{\rm root}(\theta)
  =\Tr\Bigl(M\,(m-1)\!\int\theta\theta^T\dd\rho_{\rm root}\Bigr)
  =\Tr M\le\eps,
\]
using the frame identity $(m-1)\int\theta\theta^T\dd\rho_{\rm root}=I_{H_0}$ and
$M$ symmetric on $H_0$.  On the other hand, by
Theorem~\ref{thm:sol-fiber-root-degree-two},
$K(c_r,c_r)\ge\lambda^*G(c_r,c_r)$ for every $r$, so
\[
  \eps\ \ge\ \sum_rw_r\,K(c_r,c_r)
  \ \ge\ \lambda^*\sum_rw_r\,G(c_r,c_r)=\lambda^*
  =\frac{(m+2)(m+3)}{5m^2}>\frac15 .
\]

\emph{(2).}  Immediate from (1): every degree-$2$ certificate objective exceeds
$\tfrac15$ uniformly in $m\ge3$, so no sequence $\eps_m\to0$ exists at $k=2$.
Within the candidate framework, in which a fixed-degree dual refuter is exactly
such a certificate sequence, any polynomial dual refutation must therefore use
degree $k\ge3$.

\emph{(3).}  For any admissible $\rho$ whose averaged form $A_\rho$ is defined on
$V_{m,2}$, the supremum defining $\Lambda_{m,2}$ dominates the value at
$\rho=\rho_{\rm root}$, and by \eqref{eq:sol-frd2-root-identification} and
Theorem~\ref{thm:sol-fiber-root-degree-two} that value is
$\lmin(K,G)|_{V_{m,2}}=\lambda^*$.
\end{proof}

\subsection*{Remarks, non-claims, and audit}

\begin{remark}[agreement with the recorded exact rational artifacts]
\label{rem:sol-frd2-artifacts}
The provenance-stamped run artifacts
\texttt{research/runs/2026-08-30T172853.262299Z-fiber-frame-dual.jsonl} and
\texttt{research/runs/2026-08-30T173244.394413Z-fiber-frame-dual.jsonl} record, in
exact rational arithmetic, the degree-two root-frame pencil minimum at
$m=3,4,5$ as $2/3$, $21/40$, $56/125$ respectively, and record the radial
quotient anchor $(m+2)(m+3)/(5m^2)$ exactly for every computed $m\le15$.  These
values agree with the closed form \eqref{eq:sol-frd2-main}:
$(5\cdot6)/45=2/3$, $(6\cdot7)/80=21/40$, $(7\cdot8)/125=56/125$.  Per repository
constraint, these artifacts are directional research evidence only; no step of the
proofs above uses them, and this agreement certifies nothing.
\end{remark}

\begin{remark}[unclosed steps]
\label{rem:sol-frd2-unclosed}
None.  Theorem~\ref{thm:sol-fiber-root-degree-two} is unconditional; every step is
finite-dimensional linear algebra, elementary representation theory of $S_m$
(proved from facts (F1)--(F5), themselves proved or classical with one-line
arguments), and exact Dirichlet moments
(Lemma~\ref{lem:sol-frd2-moments}, whose proof invokes the classical
Gamma--Dirichlet factorization: the normalized vector of i.i.d.\ standard
exponentials is $\Dir(1,\dots,1)$ and independent of their sum).
Corollary~\ref{cor:sol-frd2-no-degree-two-refuter} is conditional exactly on the
hypothesis stated in its preamble, namely the certified identification
\eqref{eq:sol-frd2-root-identification} (nodes \emph{lem:conditional-fiber-form}
and \emph{prop:conditional-fiber-root-obstruction}, both \texttt{checked\_by:
agent} with a persisted review), together with the candidate status of the
$\Lambda_{m,k}$/dual-certificate framework itself.
\end{remark}

\begin{remark}[explicit non-claims]
\label{rem:sol-frd2-nonclaims}
This dossier does \emph{not} claim: any upper bound on $\Lambda_{m,2}$ or on any
all-frame quantity; anything about $\Lambda_{m,k}$ or root-frame pencils for
$k\ge3$; anything about non-polynomial tests --- indeed the certified vertex-cap
obstruction (Proposition~\ref{prop:conditional-fiber-root-obstruction}) shows the
\emph{full $L^2$} root-frame gap is $O(m^{-2})$, so the degree-two floor
\eqref{eq:sol-frd2-main} genuinely does not extend beyond the polynomial quotient;
any optimality of the root orbit among admissible frames; any answer to
Question~\ref{q:conditional-fiber-frame}; and nothing about KLS.
\end{remark}

\paragraph{Obstructions respected.}
The candidate node carries no \texttt{bounded\_by} edge.  Consistency checks
against the neighboring certified and imported obstructions:
(i)~\emph{prop:conditional-fiber-root-obstruction} (root-frame $L^2$ gap
$O(m^{-2})$ via a vertex-cap indicator) is compatible with
\eqref{eq:sol-frd2-main} because the cap indicator is not a polynomial of degree
two; the two statements jointly prove that any test exhibiting the root-frame
collapse must leave $V_{m,2}$, which is precisely the content of
Corollary~\ref{cor:sol-frd2-no-degree-two-refuter}(2) for the dual channel.
(ii)~\emph{prop:sasada-negative-exchange} (imported negative-rate exchange upper
bound $48/[m(m+1)]$ for the $L^2$ gap) is compatible for the same reason.
(iii)~The six route-level obstruction nodes (\emph{obs:two-tail},
\emph{obs:proj-ceiling}, \emph{obs:crude-insufficient},
\emph{obs:relative-ceiling}, \emph{obs:circularity},
\emph{obs:rank-one-refuted}) fence Eldan-localization proof shapes; no stochastic
localization, projection summation, bootstrap, isoperimetric localization, or
rank-one inference occurs here.  No numerical output justifies any step.

% No \printbibliography here (matches modules/*.tex): main.tex prints the single
% shared bibliography.  Standalone, \ref of manuscript labels shows "??" (expected).
\end{document}
